\section*{Capítulo \ref{ch:b3:advanced-nmr} --- RMN em profundidade: pulsos, relaxação e 2D}

\begin{solution}{exo:b3:advanced-nmr:1}
$\nu_0 = (\gamma/2\pi)B_0$: \ce{^1H} \qty{600.4}{MHz}, \ce{^{13}C} \qty{151.0}{MHz}, \ce{^{19}F}
\qty{565.1}{MHz}.
\end{solution}

\begin{solution}{exo:b3:advanced-nmr:2}
$h\nu_0/2kT = 6.626\times10^{-34} \times 400.2\times10^6/(2 \times 1.381\times10^{-23} \times 300)
= \num{3.2e-5}$. A \qty{4}{K}, ele é 75 vezes maior, \num{2.4e-3} (a aproximação
ainda vale).
\end{solution}

\begin{solution}{exo:b3:advanced-nmr:3}
$\theta = \gamma B_1t_p = \pi/2$ em \qty{10}{\micro s}: $\gamma B_1/2\pi = 1/(4 \times
\qty{10}{\micro s}) = \qty{25}{kHz}$. Um pulso de $180^\circ$ dura \qty{20}{\micro s}; \qty{5}{\micro s}
dão $45^\circ$.
\end{solution}

\begin{solution}{exo:b3:advanced-nmr:4}
$T_2 = 1/(\pi \times 3.2) = \qty{0.10}{s}$.
\end{solution}

\begin{solution}{exo:b3:advanced-nmr:5}
$(10.71/42.58)^3 = 0.0159$, vezes 0.011: \num{1.75e-4}, cerca de 1/5700 do sinal
do próton. Como S/R $\propto\sqrt N$, a mesma S/R exige $5700^2 \approx \num{3e7}$ vezes mais
varreduras --- impossível, e é por isso que os espectros de ${}^{13}$C usam amostras
mais concentradas, desacoplamento e transferência de polarização.
\end{solution}

\begin{solution}{exo:b3:advanced-nmr:6}
$T_1 = 1.39/\ln2 = \qty{2.0}{s}$; tempo de espera $\ge 5T_1 = \qty{10}{s}$.
\end{solution}

\begin{solution}{exo:b3:advanced-nmr:7}
\omterm{def:b3:advanced-nmr:experiments}{COSY}: um único \omterm{def:b3:advanced-nmr:two-d}{pico cruzado}, entre o quarteto do \ce{CH2} e o tripleto do \ce{CH3} do
grupo etila; o singleto do \ce{CH3CO} não se acopla a nada. \omterm{def:b3:advanced-nmr:experiments}{HMBC} a partir do metila
singleto: o carbono da carbonila (duas ligações) e o carbono do \ce{CH2} (três ligações).
\end{solution}

\begin{solution}{exo:b3:advanced-nmr:8}
$k_c = \pi \times 50/\sqrt2 = \qty{111}{s^{-1}}$. $\Delta G^\ddagger = 8.314 \times 330 \times
\ln(6.88\times10^{12}/111) = \qty{68}{kJ/mol}$.
\end{solution}

\begin{solution}{exo:b3:advanced-nmr:9}
NOE $\propto r^{-6}$: $r_b = 0.30 \times 4^{-1/6} = \qty{0.24}{nm}$.
\end{solution}

\begin{solution}{exo:b3:advanced-nmr:10}
No \omterm{def:b3:advanced-nmr:magnetisation}{referencial girante}, em ressonância, só resta $\mathbf B_1 = B_1\mathbf e_{x'}$:
$\dot M_{x'} = 0$, $\dot M_{y'} = \gamma B_1M_z$, $\dot M_z = -\gamma B_1M_{y'}$ (componentes de
$\gamma\mathbf M\times\mathbf B_1$). Logo $M_{x'}$ é constante e $(M_{y'}, M_z)$ gira com a velocidade
angular $\gamma B_1$: depois de $t_p$, do ângulo $\gamma B_1t_p$ em torno de $x'$.
\end{solution}

\begin{solution}{exo:b3:advanced-nmr:11}
Um spin em um campo ligeiramente mais intenso ganha fase a uma taxa extra
constante durante $\tau$; o pulso de $180^\circ$ inverte sua fase, e ele perde exatamente a
mesma quantidade durante o segundo $\tau$: todos esses spins se realinham em $2\tau$. As
interações spin--spin aleatórias flutuam no tempo e não se repetem no segundo
intervalo: a defasagem que causam não é desfeita. A amplitude do eco decai com o verdadeiro $T_2$.
\end{solution}

\begin{solution}{exo:b3:advanced-nmr:12}
Etanoato de etila, \ce{CH3COOCH2CH3}: quarteto do \ce{OCH2} (4.12), singleto do \ce{CH3CO} (2.05),
tripleto do \ce{CH3}; os prótons do quarteto veem o carbono da carbonila a três ligações,
através do oxigênio do éster. O propanoato de metila, \ce{CH3CH2COOCH3}, mostraria seu
quarteto perto de 2.3 ppm (sobre um carbono perto de 27 ppm no \omterm{def:b3:advanced-nmr:experiments}{HSQC}, e não de 60) e um singleto
de \ce{OCH3} perto de 3.7 ppm, cujo pico \omterm{def:b3:advanced-nmr:experiments}{HMBC} até a carbonila atravessa o oxigênio.
\end{solution}

\begin{solution}{pb:b3:advanced-nmr:1}
\textbf{1.} $(2 \times 6 + 2 - 12)/2 = 1$: uma C=O.
\textbf{2.} 173.6 ppm: a carbonila do éster; 60.1 ppm: o \ce{OCH2}.
\textbf{3.} Cinco; DEPT-135: \ce{CH2} negativos (60.1, 36.3, 18.6), \ce{CH3} positivos (14.3,
13.7), a carbonila ausente.
\textbf{4.} \qty{400.2}{MHz} e \qty{100.7}{MHz}.
\textbf{5.} $3J = \qty{21.6}{Hz} = \qty{0.054}{ppm}$.
\textbf{6.} As multiplicidades mostram quais grupos são vizinhos dentro de cada
fragmento, mas não como os fragmentos se unem através do oxigênio e da carbonila.
\textbf{7.} 4.13--1.25; 2.28--1.66; 1.66--0.95 (e os simétricos).
\textbf{8.} \ce{OCH2CH3} (4.13, 1.25) e \ce{CH2CH2CH3} (2.28, 1.66, 0.95).
\textbf{9.} Os prótons do \ce{OCH2} e do \ce{CH2C=O} estão separados pelo oxigênio e pelo
carbono da carbonila: cinco ligações, nenhum acoplamento resolvido.
\textbf{10.} 1.66 ppm, entre \ce{CH2} e \ce{CH3}: acoplado a 2 + 3 = 5 prótons com
$J$ semelhantes, um sexteto (um multipleto).
\textbf{11.} Metilas acoplados, logo dois \ce{CH3} em carbonos adjacentes (uma unidade \ce{CH3-CH3}): impossível aqui.
\textbf{12.} \ce{-O-CH2-CH3} e \ce{-CH2-CH2-CH3}, mais o \ce{C=O}.
\textbf{13.} 4.13/60.1; 2.28/36.3; 1.66/18.6; 1.25/14.3; 0.95/13.7.
\textbf{14.} 14.3 (\ce{CH3} da etila) e 13.7 ppm (\ce{CH3} do butanoíla), a apenas \qty{0.6}{ppm} um do outro.
\textbf{15.} 4.13 ppm a 173.6 ppm: H--C--O--C, três ligações.
\textbf{16.} 2.28 ppm (duas ligações) e 1.66 ppm (três ligações) a 173.6 ppm.
\textbf{17.} \ce{CH3CH2CH2C(=O)OCH2CH3}, o butanoato de etila. O propanoato de propila poria um
tripleto de \ce{OCH2} (e não um quarteto) perto de 4.0 ppm e um quarteto perto de 2.3 ppm.
\textbf{18.} Os acoplamentos a quatro ligações costumam ser inferiores a 1 Hz; o atraso do \omterm{def:b3:advanced-nmr:experiments}{HMBC},
ajustado para 5--10 Hz, não os seleciona.
\textbf{19.} Os carbonos relaxam com velocidades diferentes e recebem do desacoplamento
intensificações Overhauser nucleares diferentes.
\textbf{20.} $1 - \eu^{-5/20} = 0.22$: a carbonila fica sub-representada por um fator 4.
\textbf{21.} $1 - \eu^{-5} = 0.993$.
\textbf{22.} O NOE do desacoplamento de prótons, diferente para cada carbono; ele é removido
pelo desacoplamento inverso controlado (desacoplador ligado apenas durante a aquisição).
\textbf{23.} $256 \times \qty{100}{s} = \qty{25600}{s}$, cerca de 7 horas.
\textbf{24.} \textbf{$5T_1$ do carbono da carbonila: \qty{100}{s} entre as varreduras.}
\end{solution}
